\section{Trust model and problem formulation}
\label{sec:model}

Let a captured form image be evidence $e\in E$, and let $M$ denote the set of machine modules exposed through the application: perception, optional language-model review, and retrieval. A module may derive a candidate $c\in C$ from evidence or existing non-authoritative context. Every persisted candidate carries a \emph{certificate}: a content address computed over its field key, value payload, evidence hash and locator, producer and template identity, selection artifact, target record, expected fact version, and creation time. The certificate is an application-level content-addressed authorization artifact, not a PKI certificate or digital identity credential. A human decision $h\in H$ records an actor, a per-field action, an expected record version, a reason, and the referenced certificate. Authoritative state is represented as a sequence of fact versions $F_0,F_1,\ldots$.

The allowed transitions are
\begin{equation}
  E\xrightarrow{m}C,\qquad C\xrightarrow{\text{review}}H,
  \qquad F_{t+1}=g(F_t,H_t),
  \label{eq:transitions}
\end{equation}
where $g$ validates the expected version and appends, rather than overwrites, the resulting record. The transition is per field: one attributable decision and one fact transition per field, bound to the exact certificate the human relied on, with the certificate's field key, evidence binding, and expected version checked before any write, and the whole version committed by one compare-and-swap on the current fact version. Retrieval can add a reference to $C$ but cannot supply $H$. The central invariant is
\begin{equation}
  \forall m\in M,\quad m\not\rightarrow F.
  \label{eq:invariant}
\end{equation}
Equation~\ref{eq:invariant} is an application-architecture property: no exposed machine service possesses a direct fact-writing operation. It is not a claim about protection against privileged database modification, compromised dependencies, or an adversary controlling the host.

In this prototype, human authorization denotes an attributable review action issued through the review capability; reviewer identity is recorded but not cryptographically authenticated. The evaluated invariant concerns capability separation rather than proof of human presence. Figure~\ref{fig:trust-model} shows the trust model.

\begin{figure}[t]
  \centering
  \resizebox{0.96\textwidth}{!}{\begin{tikzpicture}[node distance=7mm and 8mm,font=\small]
  \node[machine,minimum width=29mm] (perception) {Perception};
  \node[machine,minimum width=29mm,below=of perception] (llm) {Optional LLM};
  \node[machine,minimum width=29mm,below=of llm] (retrieval) {Retrieval};
  \node[trustbox,minimum width=34mm,right=20mm of llm] (candidate) {Candidate plane\\certificates, append-only};
  \node[human,minimum width=34mm,right=17mm of candidate] (decision) {Human decision\\confirm/correct};
  \node[fact,minimum width=34mm,right=17mm of decision] (fact) {Fact plane\\versioned};
  \draw[flow] (perception.east) -- (candidate.west);
  \draw[flow] (llm.east) -- (candidate.west);
  \draw[flow] (retrieval.east) -- node[below,sloped,font=\scriptsize]{read-only reference} (candidate.west);
  \draw[flow] (candidate) -- (decision);
  \draw[flow] (decision) -- (fact);
  \draw[noflow] ([yshift=-10mm]llm.south east) -- node[below,font=\scriptsize,text=TrustAmber]{no machine write path} ([yshift=-10mm]fact.south west);
  \node[draw=TrustCharcoal,dotted,rounded corners,fit=(perception)(retrieval),inner sep=4mm,label={[text=TrustCharcoal]above:untrusted machine modules}] {};
\end{tikzpicture}}
  \caption{Trust model. Dashed machine components terminate at append-only candidates or read-only references; an attributable human decision is the only exposed path to a versioned fact.}
  \label{fig:trust-model}
\end{figure}

For a recognition input $x$, a predictor returns a class $\hat y(x)$ and a selection score $s(x)$. Given threshold $\tau$, the selection function is
\begin{equation}
  a_\tau(x)=\mathbb{1}[s(x)\geq\tau].
\end{equation}
For a sample of $n$ labeled cases, coverage and selective risk are
\begin{align}
  \widehat{\mathrm{cov}}(\tau)&=\frac{1}{n}\sum_{i=1}^{n}a_\tau(x_i),\\
  \widehat{R}(\tau)&=\frac{\sum_{i=1}^{n}a_\tau(x_i)\mathbb{1}[\hat y_i\neq y_i]}{\max(1,\sum_{i=1}^{n}a_\tau(x_i))}.
\end{align}
The erroneous auto-pass rate uses all cases in the denominator. These empirical measures characterize the recognition gate; they do not alter the authority rule. An accepted recognition is still a candidate, not a fact.

The desired properties are: (P1) machine isolation, Equation~\ref{eq:invariant}; (P2) append-only candidate and fact histories whose candidates are content-addressed certificates; (P3) optimistic concurrency, so a stale decision cannot overwrite a newer version; (P4) content-addressed evidence and reverse trace from export to its supporting chain, rebuilt per field as evidence--certificate--decision--transition with any tampered or missing binding reported as incomplete; and (P5) graceful operation when the optional AI adapter is disabled.
